\documentclass[10pt,a4paper]{amsart}
\usepackage{amsmath,amssymb,amsthm,graphicx}
% Retained preamble for the Stacks Project chapter snapshot in this directory.
% The chapter's own preamble.tex is a copy of the shared Stacks Project preamble; the block
% below is sliced MECHANICALLY, line for line, from it. Nothing was added, removed or
% redefined; the wrapper class, the xr-hyper/hyperref wiring and the externaldocument list
% are omitted because the assembler supplies the document class.
%
% lines taken from preamble.tex: 3, 5, 9, 11, 12, 156, 157, 158, ... (51 lines total)
%
% package coverage restored from the chapter's own preamble (2026-10-13): xy, xy-2cell

\usepackage[all]{xy}
\xyoption{2cell}
\UseAllTwocells

\newenvironment{reference}{\comment}{\endcomment}
\newenvironment{slogan}{\comment}{\endcomment}
\newenvironment{history}{\comment}{\endcomment}
\theoremstyle{plain}
\newtheorem{theorem}[subsection]{Theorem}
\newtheorem{proposition}[subsection]{Proposition}
\newtheorem{lemma}[subsection]{Lemma}
\theoremstyle{definition}
\newtheorem{definition}[subsection]{Definition}
\newtheorem{example}[subsection]{Example}
\newtheorem{exercise}[subsection]{Exercise}
\newtheorem{situation}[subsection]{Situation}
\theoremstyle{remark}
\newtheorem{remark}[subsection]{Remark}
\newtheorem{remarks}[subsection]{Remarks}
\numberwithin{equation}{subsection}
\def\lim{\mathop{\mathrm{lim}}\nolimits}
\def\colim{\mathop{\mathrm{colim}}\nolimits}
\def\Spec{\mathop{\mathrm{Spec}}}
\def\Hom{\mathop{\mathrm{Hom}}\nolimits}
\def\Ext{\mathop{\mathrm{Ext}}\nolimits}
\def\SheafHom{\mathop{\mathcal{H}\!\mathit{om}}\nolimits}
\def\SheafExt{\mathop{\mathcal{E}\!\mathit{xt}}\nolimits}
\def\Sch{\mathit{Sch}}
\def\Mor{\mathop{\mathrm{Mor}}\nolimits}
\def\Ob{\mathop{\mathrm{Ob}}\nolimits}
\def\Sh{\mathop{\mathit{Sh}}\nolimits}
\def\NL{\mathop{N\!L}\nolimits}
\def\CH{\mathop{\mathrm{CH}}\nolimits}
\def\proetale{{pro\text{-}\acute{e}tale}}
\def\etale{{\acute{e}tale}}
\def\QCoh{\mathit{QCoh}}
\def\Ker{\mathop{\mathrm{Ker}}}
\def\Im{\mathop{\mathrm{Im}}}
\def\Coker{\mathop{\mathrm{Coker}}}
\def\Coim{\mathop{\mathrm{Coim}}}
\DeclareMathSymbol{\boxtimes}{\mathbin}{AMSa}{"02}
\def\QCohstack{\mathcal{QC}\!\mathit{oh}}
\def\Cohstack{\mathcal{C}\!\mathit{oh}}
\def\Spacesstack{\mathcal{S}\!\mathit{paces}}
\def\Quotfunctor{\mathrm{Quot}}
\def\Hilbfunctor{\mathrm{Hilb}}
\def\Curvesstack{\mathcal{C}\!\mathit{urves}}
\def\Polarizedstack{\mathcal{P}\!\mathit{olarized}}
\def\Complexesstack{\mathcal{C}\!\mathit{omplexes}}
\def\Pic{\mathop{\mathrm{Pic}}\nolimits}
\def\Picardstack{\mathcal{P}\!\mathit{ic}}
\def\Picardfunctor{\mathrm{Pic}}
\def\Deformationcategory{\mathcal{D}\!\mathit{ef}}
\usepackage{cleveref}
\crefname{originallemma}{original lemma}{original lemmas}
\Crefname{originallemma}{Original lemma}{Original lemmas}
\makeatletter
\newlabel{lemma-points-cartesian}{{lemma-points-cartesian}{source}{}{sourcepointer-1}{}}
\newlabel{lemma-points-cartesian@cref}{{[originallemma][0][]lemma-points-cartesian}{[1][0][]source}{}{}{}}
\makeatother
\crefname{originallemma}{original lemma}{original lemmas}
\Crefname{originallemma}{Original lemma}{Original lemmas}
\makeatletter
\newlabel{section-properties-morphisms}{{section-properties-morphisms}{source}{}{sourcepointer-2}{}}
\newlabel{section-properties-morphisms@cref}{{[originallemma][0][]section-properties-morphisms}{[1][0][]source}{}{}{}}
\makeatother
\crefname{originallemma}{original lemma}{original lemmas}
\Crefname{originallemma}{Original lemma}{Original lemmas}
\makeatletter
\newlabel{spaces-morphisms-definition-surjective}{{03ME}{source}{}{sourcepointer-3}{}}
\newlabel{spaces-morphisms-definition-surjective@cref}{{[originallemma][0][]03ME}{[1][0][]source}{}{}{}}
\makeatother
\title{For representable morphisms, point-surjectivity and section-surjectivity agree}
\author{The Stacks Project Authors}
\date{stacks-project:stacks-properties.tex:local-snapshot-20261012; GNU FDL 1.2 or later}
\begin{document}
\maketitle
\noindent\textit{Editorial note.} Statement span is lines 1050-1059 of sources/stacks-stacks-properties/pinned-original-source.tex; proof span is lines 1061-1095, the complete author proof (33 body lines).\par\medskip
\noindent\textit{Editorial note.} Navigation stubs: 3 cited-but-not-included labels are supplied as non-load-bearing pointers with their pinned tags where the snapshot tag index has them.\par\medskip
\noindent\textit{Editorial note.} extra\_packages is empty on purpose: the retained preamble already loads xy with options and a second declaration clashes in the assembler wrapper.\par\medskip
\noindent\textit{Editorial note.} Text is retained verbatim, including the author's own typos and placeholders.\par\medskip
\noindent\textit{Editorial note.} Queue index 16 of queue\_stacks-properties.json (pb=1204, stmt 1050-1059, proof 1061-1095).\par\medskip
\noindent\textit{Editorial note.} stacks-properties wave-1: fresh adjudication accept (2026-10-14); pool scanned while CONTROL.md is down (0 bytes) - §2 row deferred per Chief, claim reported to Chief 08:0x.\par\medskip
\noindent\textit{Allowed standard prerequisites (editorial).}\begin{itemize}
\item Explicit prerequisite boundary. The cited inputs are prerequisites; the argument in the retained span is not replaced by any cited result.
\end{itemize}\medskip
\ifdefined\hypertarget\hypertarget{sourcepointer-1}{}\fi
\noindent\textit{Source reference.} Original lemma lemma-points-cartesian: sources/stacks-stacks-properties/pinned-original-source.tex, source lines 1019--1036. Cited by the retained statement or proof but not part of the collected spans; supplied only as a navigation pointer. The argument does not depend on its content.\par\medskip
\ifdefined\hypertarget\hypertarget{sourcepointer-2}{}\fi
\noindent\textit{Source reference.} Original lemma section-properties-morphisms: sources/stacks-stacks-properties/pinned-original-source.tex, source lines 378--401. Cited by the retained statement or proof but not part of the collected spans; supplied only as a navigation pointer. The argument does not depend on its content.\par\medskip
\ifdefined\hypertarget\hypertarget{sourcepointer-3}{}\fi
\noindent\textit{Source reference.} Original lemma 03ME: sources/stacks-stacks-properties/pinned-original-source.tex, source lines 1095--1095. Cited by the retained statement or proof but not part of the collected spans; supplied only as a navigation pointer. The argument does not depend on its content.\par\medskip
\begin{lemma}
\label{lemma-characterize-surjective}
Let $f : \mathcal{X} \to \mathcal{Y}$ be a morphism of algebraic stacks
which is representable by algebraic spaces. The following are equivalent:
\begin{enumerate}
\item $|f| : |\mathcal{X}| \to |\mathcal{Y}|$ is surjective, and
\item $f$ is surjective
(in the sense of Section \ref{section-properties-morphisms}).
\end{enumerate}
\end{lemma}

\begin{proof}
Assume (1). Let $T \to \mathcal{Y}$ be a morphism whose source is a scheme.
To prove (2) we have to show that the morphism of algebraic spaces
$T \times_\mathcal{Y} \mathcal{X} \to T$ is surjective. By
Morphisms of Spaces, Definition \ref{spaces-morphisms-definition-surjective}
this means we have to show that
$|T \times_\mathcal{Y} \mathcal{X}| \to |T|$ is surjective.
Applying
Lemma \ref{lemma-points-cartesian}
we see that this follows from (1).

\medskip\noindent
Conversely, assume (2). Let $y : \Spec(K) \to \mathcal{Y}$ be a
morphism from the spectrum of a field into $\mathcal{Y}$. By assumption the
morphism
$\Spec(K) \times_{y, \mathcal{Y}} \mathcal{X} \to \Spec(K)$
of algebraic spaces is surjective. By
Morphisms of Spaces, Definition \ref{spaces-morphisms-definition-surjective}
this means there exists a field extension
$K'/K$ and a morphism
$\Spec(K') \to \Spec(K) \times_{y, \mathcal{Y}} \mathcal{X}$
such that the left square of the diagram
$$
\xymatrix{
\Spec(K') \ar[r] \ar[d] &
\Spec(K) \times_{y, \mathcal{Y}} \mathcal{X} \ar[d] \ar[r] &
\mathcal{X} \ar[d]
\\
\Spec(K) \ar@{=}[r] &
\Spec(K) \ar[r]^-y &
\mathcal{Y}
}
$$
is commutative. This shows that $|X| \to |\mathcal{Y}|$ is surjective.
\end{proof}

\end{document}
